\documentclass[11pt,compress,t,notes=noshow, xcolor=table, aspectratio=169]{beamer}
\usepackage[]{graphicx}

\usepackage[]{color}
% maxwidth is the original width if it is less than linewidth
% otherwise use linewidth (to make sure the graphics do not exceed the margin)
\makeatletter
\def\maxwidth{ %
  \ifdim\Gin@nat@width>\linewidth
    \linewidth
  \else
    \Gin@nat@width
  \fi
}
\makeatother

\definecolor{fgcolor}{rgb}{0.345, 0.345, 0.345}
\newcommand{\hlnum}[1]{\textcolor[rgb]{0.686,0.059,0.569}{#1}}%
\newcommand{\hlstr}[1]{\textcolor[rgb]{0.192,0.494,0.8}{#1}}%
\newcommand{\hlcom}[1]{\textcolor[rgb]{0.678,0.584,0.686}{\textit{#1}}}%
\newcommand{\hlopt}[1]{\textcolor[rgb]{0,0,0}{#1}}%
\newcommand{\hlstd}[1]{\textcolor[rgb]{0.345,0.345,0.345}{#1}}%
\newcommand{\hlkwa}[1]{\textcolor[rgb]{0.161,0.373,0.58}{\textbf{#1}}}%
\newcommand{\hlkwb}[1]{\textcolor[rgb]{0.69,0.353,0.396}{#1}}%
\newcommand{\hlkwc}[1]{\textcolor[rgb]{0.333,0.667,0.333}{#1}}%
\newcommand{\hlkwd}[1]{\textcolor[rgb]{0.737,0.353,0.396}{\textbf{#1}}}%
\let\hlipl\hlkwb

\usepackage{framed}
\makeatletter
\newenvironment{kframe}{%
 \def\at@end@of@kframe{}%
 \ifinner\ifhmode%
  \def\at@end@of@kframe{\end{minipage}}%
  \begin{minipage}{\columnwidth}%
 \fi\fi%
 \def\FrameCommand##1{\hskip\@totalleftmargin \hskip-\fboxsep
 \colorbox{shadecolor}{##1}\hskip-\fboxsep
     % There is no \\@totalrightmargin, so:
     \hskip-\linewidth \hskip-\@totalleftmargin \hskip\columnwidth}%
 \MakeFramed {\advance\hsize-\width
   \@totalleftmargin\z@ \linewidth\hsize
   \@setminipage}}%
 {\par\unskip\endMakeFramed%
 \at@end@of@kframe}
\makeatother

\definecolor{shadecolor}{rgb}{.97, .97, .97}
\definecolor{messagecolor}{rgb}{0, 0, 0}
\definecolor{warningcolor}{rgb}{1, 0, 1}
\definecolor{errorcolor}{rgb}{1, 0, 0}
\newenvironment{knitrout}{}{} % an empty environment to be redefined in TeX

\newcommand{\SweaveOpts}[1]{}  % do not interfere with LaTeX
\newcommand{\SweaveInput}[1]{} % because they are not real TeX commands
\newcommand{\Sexpr}[1]{}       % will only be parsed by R
\newcommand{\xmark}{\ding{55}}%

\usepackage{alltt}
\usepackage[utf8]{inputenc}

\usepackage[english]{babel}
\usepackage{dsfont}
\usepackage{verbatim}
\usepackage{amsmath}
\usepackage{amsfonts}
\usepackage{amssymb}
\usepackage{bm}
\usepackage{csquotes}
\usepackage{multirow}
\usepackage{longtable}
\usepackage{booktabs}
\usepackage{enumerate}
\usepackage[absolute,overlay]{textpos}
\usepackage{psfrag}
\usepackage{algorithm}
\usepackage{algpseudocode}
\usepackage{eqnarray}
\usepackage{arydshln}
\usepackage{tabularx}
\usepackage{placeins}
\usepackage{tikz}
\usepackage{setspace}
\usepackage{colortbl}
\usepackage{mathtools}
\usepackage{wrapfig}
\usepackage{bm}
\usepackage{amsmath}
\usepackage{pifont}
\usepackage{xcolor} %colored math symbols
\usepackage{adjustbox}
\usepackage{verbatimbox}
\usepackage{tcolorbox}
\usepackage{siunitx}
\def\signed #1{{\leavevmode\unskip\nobreak\hfil\penalty50\hskip1em
  \hbox{}\nobreak\hfill #1%
  \parfillskip=0pt \finalhyphendemerits=0 \endgraf}}

% basic latex stuff
\newcommand{\col}{\par\colorbox{code}{\parbox{\textwidth}{\theverbbox}}\par}


\usetikzlibrary{shapes,arrows,automata,positioning,calc,chains,trees, shadows}
\tikzset{
  %Define standard arrow tip
  >=stealth',
  %Define style for boxes
  punkt/.style={
    rectangle,
    rounded corners,
    draw=black, very thick,
    text width=6.5em,
    minimum height=2em,
    text centered},
  % Define arrow style
  pil/.style={
    ->,
    thick,
    shorten <=2pt,
    shorten >=2pt,}
}

\newsavebox\mybox
\newenvironment{aquote}[1]
  {\savebox\mybox{#1}\begin{quote}\openautoquote\hspace*{-.7ex}}
  {\unskip\closeautoquote\vspace*{1mm}\signed{\usebox\mybox}\end{quote}}

\usepackage{subfig}

% Defines macros and environments
\usepackage{../style/dhbw-lecture}


\let\code=\texttt
\let\proglang=\textsf

\setkeys{Gin}{width=0.9\textwidth}

% The accent rule sits on its own vbox line with \nointerlineskip so it costs
% only its own 1.2pt height + the explicit 2pt gap, instead of pulling in a
% full \baselineskip (~0.6-0.7cm) the way a plain "\\" line break would --
% that overhead on every single frame is what was pushing picture-heavy
% frames over the shorter 16:9 frame height.
\setbeamertemplate{frametitle}{\expandafter\uppercase\expandafter\insertframetitle\par\vskip2pt\nointerlineskip\hbox{\color{dhbwaccent}\rule{1.6cm}{1.2pt}}}
\addtobeamertemplate{frametitle}{}{\vspace{0.08cm}}

\usepackage{bbm}
% basic latex stuff
\newcommand{\pkg}[1]{{\fontseries{b}\selectfont #1}} %fontstyle for R packages
\newcommand{\lz}{\vspace{0.5cm}} %vertical space
\newcommand{\dlz}{\vspace{1cm}} %double vertical space
\newcommand{\oneliner}[1] % Oneliner for important statements
{\begin{block}{}\begin{center}\begin{Large}#1\end{Large}\end{center}\end{block}}


%new environments
\newenvironment{vbframe}  %frame with breaks and verbatim
{
 \begin{frame}[containsverbatim,allowframebreaks]
}
{
\end{frame}
}

\newenvironment{vframe}  %frame with verbatim without breaks (to avoid numbering one slided frames)
{
 \begin{frame}[containsverbatim]
}
{
\end{frame}
}

\newenvironment{blocki}[1]   % itemize block
{
 \begin{block}{#1}\begin{itemize}
}
{
\end{itemize}\end{block}
}

\newenvironment{fragileframe}[2]{  %fragile frame with framebreaks
\begin{frame}[allowframebreaks, fragile, environment = fragileframe]
\frametitle{#1}
#2}
{\end{frame}}


\newcommand{\myframe}[2]{  %short for frame with framebreaks
\begin{frame}[allowframebreaks]
\frametitle{#1}
#2
\end{frame}}

\newcommand{\remark}[1]{
  \textbf{Remark:} #1
}


\newenvironment{deleteframe}
{
\begingroup
\usebackgroundtemplate{\includegraphics[width=\paperwidth,height=\paperheight]{../style/color/red.png}}
 \begin{frame}
}
{
\end{frame}
\endgroup
}
\newenvironment{simplifyframe}
{
\begingroup
\usebackgroundtemplate{\includegraphics[width=\paperwidth,height=\paperheight]{../style/color/yellow.png}}
 \begin{frame}
}
{
\end{frame}
\endgroup
}\newenvironment{draftframe}
{
\begingroup
\usebackgroundtemplate{\includegraphics[width=\paperwidth,height=\paperheight]{../style/color/green.jpg}}
 \begin{frame}
}
{
\end{frame}
\endgroup
}
% https://tex.stackexchange.com/a/261480: textcolor that works in mathmode
\makeatletter
\renewcommand*{\@textcolor}[3]{%
  \protect\leavevmode
  \begingroup
    \color#1{#2}#3%
  \endgroup
}
\makeatother


\input{../latex-math/basic-math}
\input{../latex-math/basic-ml}
\input{../latex-math/ml-nn}

% The course's typewriter font (laett) has no bold shape; map bold typewriter to regular.
\AtBeginDocument{{\ttfamily\selectfont}\DeclareFontShape{T1}{laett}{bx}{n}{<->ssub*laett/m/n}{}\DeclareFontShape{T1}{laett}{m}{bx}{<->ssub*laett/m/n}{}}
\newcommand{\citebutton}[2]{\beamergotobutton{\href{#2}{#1}}}

\newcommand{\titlefigure}{figure/title-modern-variant.png}
\newcommand{\learninggoals}{
  \item What changed in the Transformer block since 2017 and why: pre-norm, RMSNorm, no biases, gated activations
  \item Rotary position embeddings (RoPE)
  \item The consensus hyperparameters of modern language models
  \item Stability tricks and attention variants (MQA/GQA, sliding windows), and why memory matters
  \item Mixture of experts: routing, training, DeepSeek-MoE
}

\title{Advanced Artificial Intelligence and Machine Learning}
\date{}

\begin{document}

\lecturechapter{Transformers Modernized}
\lecture{Advanced AI \& ML}

%%%%%%%%%%%%%%%%%%%%%%%%%%%%%%%%%%%%%%%%%%%%%%%%%%%%%%%%%%%%%%%%%%
\section{Since 2017: What Changed Inside the Transformer?}
% Source: CS336 Spring 2025, lecture 3 (architecture), slides 3-8
%%%%%%%%%%%%%%%%%%%%%%%%%%%%%%%%%%%%%%%%%%%%%%%%%%%%%%%%%%%%%%%%%%

% [O] original bridge frame (draft for the user to rewrite); listed in order.txt
\begin{frame}{Since 2017: What Changed Inside the Transformer?}
\begin{itemize}
  \item Session 7 introduced the original 2017 block: post-norm with LayerNorm, sinusoidal position embeddings, a ReLU feed-forward network. GPT, BERT and T5 (session 8) already changed some of these details, e.g.\ learned or relative position embeddings.
  \item Today's large models (e.g.\ the LLaMA family, DeepSeek) keep the overall design, but change several components \emph{inside} the block.
  \item This session goes through these changes and the evidence behind them: normalization, activations, position embeddings, hyperparameters, stability, attention heads.
  \item Then a change at the level of the whole layer: \textbf{mixture of experts}.
  \item Source: Stanford CS336, Lectures 3 and 4 (Spring 2025).
\end{itemize}
\end{frame}

% FRAMES-BEGIN l3 3 8
% CS336 slide 3
\begin{frame}{Outline and goals}
\begin{itemize}
  \item Quick recap of the `standard' transformer
  \begin{itemize}
    \item What do most of the large LMs have in common?
  \end{itemize}
  \item What are common variations to the architecture / training process?
  \begin{itemize}
    \item Today's theme: the best way to learn is hands-on experience
    \item the second best way is to try to learn from others' experience
  \end{itemize}
\end{itemize}
\end{frame}

% CS336 slide 4
\begin{frame}{Starting point: the `original' transformer}
% full-body crop (figures would be too small next to the bullets)
\begin{center}
  \includegraphics[page=4,trim=20 3 8 64,clip,width=330pt]{figure/cs336-l3.pdf}
\end{center}
\end{frame}

% CS336 slide 5
\begin{frame}{A Simple, Modern Variant}
% full-body crop (figures would be too small next to the bullets)
\begin{center}
  \includegraphics[page=5,trim=20 3 8 64,clip,width=330pt]{figure/cs336-l3.pdf}
\end{center}
\end{frame}

% CS336 slide 6
\begin{frame}{How should we think about architectures?}
\begin{center}
  \includegraphics[page=6,trim=32 11 46 61,clip,width=318pt]{figure/cs336-l3.pdf}\par
\end{center}
\end{frame}

% CS336 slide 7
\begin{frame}[shrink=8]{Let's look at the data (on dense architectures)}
% full-body crop (figures would be too small next to the bullets)
\begin{center}
  \includegraphics[page=7,trim=20 3 8 64,clip,width=330pt]{figure/cs336-l3.pdf}
\end{center}
\end{frame}

% CS336 slide 8
\begin{frame}{What are we going to cover?}
\begin{itemize}
  \item Common architecture variations
  \begin{itemize}
    \item Activations, FFN
    \item Attention variants
    \item Position embeddings
  \end{itemize}
  \item Hyperparameters that (do or don't) matter
  \begin{itemize}
    \item What is ff\_dim? Do multi\_head dims always sum to model\_dim?
    \item How many vocab elements?
  \end{itemize}
  \item Stability tricks
\end{itemize}
\end{frame}

% FRAMES-END

%%%%%%%%%%%%%%%%%%%%%%%%%%%%%%%%%%%%%%%%%%%%%%%%%%%%%%%%%%%%%%%%%%
\section{Architecture Variations: Normalization}
% Source: CS336 Spring 2025, lecture 3 (architecture), slides 9-14
%%%%%%%%%%%%%%%%%%%%%%%%%%%%%%%%%%%%%%%%%%%%%%%%%%%%%%%%%%%%%%%%%%

% FRAMES-BEGIN l3 9 14
% CS336 slide 9
\begin{frame}{Architecture variations..}
\begin{itemize}
  \item Let's think about the core architecture piece
  \begin{itemize}
    \item Trends toward `LLaMA- like' architectures
  \end{itemize}
\end{itemize}
\begin{center}
  \includegraphics[page=9,trim=20 10 49 98,clip,width=265pt]{figure/cs336-l3.pdf}\par
\end{center}
\end{frame}

% CS336 slide 10
\begin{frame}{Pre-vs-post norm}
% full-body crop (figures would be too small next to the bullets)
\begin{center}
  \includegraphics[page=10,trim=20 3 8 64,clip,width=330pt]{figure/cs336-l3.pdf}
\end{center}
\end{frame}

% CS336 slide 11
\begin{frame}{Pre-vs-post-norm, the data}
% full-body crop (figures would be too small next to the bullets)
\begin{center}
  \includegraphics[page=11,trim=20 3 8 64,clip,width=330pt]{figure/cs336-l3.pdf}
\end{center}
\end{frame}

% CS336 slide 12
\begin{frame}{Pre-vs-post norm, explanations?}
% full-body crop (figures would be too small next to the bullets)
\begin{center}
  \includegraphics[page=12,trim=20 3 8 64,clip,width=330pt]{figure/cs336-l3.pdf}
\end{center}
\end{frame}

% CS336 slide 13
\begin{frame}{New things -- `double' norm.}
\begin{itemize}
  \item If putting LayerNorms in residual streams is bad.. Why not post-norm outside the stream?
  \item Recent models: Grok, Gemma 2. Olmo 2 only does non-residual post norm
\end{itemize}
\begin{center}
  \includegraphics[page=13,trim=246 56 246 118,clip,width=119pt]{figure/cs336-l3.pdf}\par
\end{center}
\end{frame}

% CS336 slide 14
\begin{frame}{LayerNorm vs RMSNorm}
% full-body crop (figures would be too small next to the bullets)
\begin{center}
  \includegraphics[page=14,trim=20 3 8 64,clip,width=330pt]{figure/cs336-l3.pdf}
\end{center}
\end{frame}

% FRAMES-END

% [O] original bridge frame (draft for the user to rewrite); listed in order.txt
\begin{frame}{Compute vs.\ Memory}
\begin{itemize}
  \item A GPU can do arithmetic much faster than it can move data between its memory and its compute units.
  \item The number of arithmetic operations per byte moved is an operation's \textbf{arithmetic intensity}. With low arithmetic intensity, an operation is limited by \textbf{memory traffic}, not by the number of floating-point operations (FLOPs).
  \item Matrix multiplications do many operations per byte loaded; normalizations and element-wise operations do few.
  \item So an operation can account for a tiny share of the FLOPs and still a noticeable share of the runtime --- the next slides use exactly this argument.
  \item Hardware details follow in session 12.
\end{itemize}
\end{frame}

% FRAMES-BEGIN l3 15 19
% CS336 slide 15
\begin{frame}{Why RMSNorm?}
\begin{itemize}
  \item Modern explanation -- it's faster (and just as good).
  \begin{itemize}
    \item Fewer operations (no mean calculation)
    \item Fewer parameters (no bias term to store)
  \end{itemize}
  \item Does this explanation make sense?
  \begin{itemize}
    \item Matrix multiplies are the vast majority of FLOPs (and memory)
  \end{itemize}
\end{itemize}
\begin{center}
  \includegraphics[page=15,trim=464 214 23 112,clip,width=95pt]{figure/cs336-l3.pdf}\par
  \includegraphics[page=15,trim=247 61 250 244,clip,width=91pt]{figure/cs336-l3.pdf}\par
\end{center}
{\scriptsize [Ivanov et al 2023]\par}
\end{frame}

% CS336 slide 16
\begin{frame}[shrink=7]{Why RMSNorm (2)}
\begin{itemize}
  \item Important lesson: FLOPS are not runtime!
  \begin{itemize}
    \item RMSNorm can still matter due to the importance of data movement
  \end{itemize}
\end{itemize}
\begin{center}
  \includegraphics[page=16,trim=470 79 143 107,clip,width=42pt]{figure/cs336-l3.pdf}\par
  \includegraphics[page=16,trim=100 185 336 130,clip,width=111pt]{figure/cs336-l3.pdf}\par
\end{center}
{\scriptsize Left top (''43G'') is FLOPS Right top (``153'') is the FLOP-to-memory ratio; [Ivanov et al 2023]\par}
\end{frame}

% CS336 slide 17
\begin{frame}{RMSNorm - validation}
\begin{itemize}
  \item RMSNorm runtime (and surprisingly, perf) gains have been seen in papers
\end{itemize}
\begin{center}
  \includegraphics[page=17,trim=59 70 27 170,clip,width=330pt]{figure/cs336-l3.pdf}\par
\end{center}
\end{frame}

% CS336 slide 18
\begin{frame}{More generally: dropping bias terms}
\begin{center}
  \includegraphics[page=18,trim=62 34 35 67,clip,width=330pt]{figure/cs336-l3.pdf}\par
\end{center}
\end{frame}

% CS336 slide 19
\begin{frame}{LayerNorm: recap}
\begin{itemize}
  \item Basically everyone does pre-norm.
  \item Intuition -- keep the good parts of residual connections
  \item Observations -- nicer gradient propagation, fewer spike
  \item Some people add a second norm outside the residual stream (NOT post-norm)
  \item Most people do RMSnorm
  \item In practice, works as well as LayerNorm
  \item But, has fewer parameters to move around, which saves on wallclock time
  \item People more generally drop bias terms since the compute/param tradeoffs are not great.
\end{itemize}
\end{frame}

% FRAMES-END

%%%%%%%%%%%%%%%%%%%%%%%%%%%%%%%%%%%%%%%%%%%%%%%%%%%%%%%%%%%%%%%%%%
\section{Architecture Variations: Activations and Layers}
% Source: CS336 Spring 2025, lecture 3 (architecture), slides 20-29
%%%%%%%%%%%%%%%%%%%%%%%%%%%%%%%%%%%%%%%%%%%%%%%%%%%%%%%%%%%%%%%%%%

% FRAMES-BEGIN l3 20 29
% CS336 slide 20
\begin{frame}{Activations}
\begin{itemize}
  \item A whole zoo of activations ..
  \item ReLU, GeLU, Swish, ELU, GLU, GeGLU, ReGLU, SeLU, SwiGLU, LiGLU
  \begin{itemize}
    \item What are these things? What do people use? Does it matter?
  \end{itemize}
\end{itemize}
\end{frame}

% CS336 slide 21
\begin{frame}{A few of the common activations}
% full-body crop (figures would be too small next to the bullets)
\begin{center}
  \includegraphics[page=21,trim=20 3 8 64,clip,width=330pt]{figure/cs336-l3.pdf}
\end{center}
\end{frame}

% CS336 slide 22
\begin{frame}{Gated activations (*GLU)}
\begin{itemize}
  \item GLUs modify the `first part' of a FF layer
  \item Instead of a linear + ReLU, augment the above with an (entrywise) linear term
  \item This gives the gated variant (ReGLU) -- note that we have an extra parameter (V)
\end{itemize}
\begin{center}
  \includegraphics[page=22,trim=169 229 184 119,clip,width=220pt]{figure/cs336-l3.pdf}\par
  \includegraphics[page=22,trim=176 115 176 233,clip,width=220pt]{figure/cs336-l3.pdf}\par
  \includegraphics[page=22,trim=176 13 178 334,clip,width=220pt]{figure/cs336-l3.pdf}\par
\end{center}
\end{frame}

% CS336 slide 23
\begin{frame}{Gated variants of standard FF layers}
% full-body crop (figures would be too small next to the bullets)
\begin{center}
  \includegraphics[page=23,trim=20 3 8 64,clip,width=330pt]{figure/cs336-l3.pdf}
\end{center}
\end{frame}

% CS336 slide 24
\begin{frame}{Do gated linear units work?}
\begin{center}
  \includegraphics[page=24,trim=72 47 135 71,clip,width=295pt]{figure/cs336-l3.pdf}\par
\end{center}
\end{frame}

% CS336 slide 25
\begin{frame}{Do gated linear units work (2)?}
\begin{itemize}
  \item Yes, with other works corroborating Shazeer 2020
  \begin{itemize}
    \item Narang et al 2020
  \end{itemize}
\end{itemize}
\begin{center}
  \includegraphics[page=25,trim=80 95 51 105,clip,width=330pt]{figure/cs336-l3.pdf}\par
\end{center}
\end{frame}

% CS336 slide 26
\begin{frame}{Gating, activations}
\begin{itemize}
  \item Many variations (ReLU, GeLU, *GLU) across models.
  \item *GLU isn't necessary for a good model (see GPT3), but it's probably helpful Otuer, recent outlier models.. Nemotron 340B (Squared ReLU), Falcon 2 11b (ReLU)
  \item But evidence points towards somewhat consistent gains from Swi/GeGLU
\end{itemize}
\end{frame}

% CS336 slide 27
\begin{frame}{Serial vs Parallel layers}
\begin{itemize}
  \item Normal transformer blocks are serial -- they compute attention, then the MLP
  \begin{itemize}
    \item Could we parallelize the transformer block?
  \end{itemize}
\end{itemize}
\begin{center}
  \includegraphics[page=27,trim=284 83 284 103,clip,width=85pt]{figure/cs336-l3.pdf}\par
\end{center}
\end{frame}

% CS336 slide 28
\begin{frame}{Parallel layers}
\begin{itemize}
  \item A few models (GPTJ, PaLM, GPT-NeoX) do parallel layers. Originally in GPT-J
  \item If implemented right, LayerNorm can be shared, and matrix multiplies can be fused
  \item Recent Models: Cohere Command A, Falcon 2 11B, Command R+
\end{itemize}
\begin{center}
  \includegraphics[page=28,trim=91 96 55 115,clip,width=310pt]{figure/cs336-l3.pdf}\par
\end{center}
\end{frame}

% CS336 slide 29
\begin{frame}{Summary: architectures}
% full-body crop (figures would be too small next to the bullets)
\begin{center}
  \includegraphics[page=29,trim=20 3 8 64,clip,width=330pt]{figure/cs336-l3.pdf}
\end{center}
\end{frame}

% FRAMES-END

%%%%%%%%%%%%%%%%%%%%%%%%%%%%%%%%%%%%%%%%%%%%%%%%%%%%%%%%%%%%%%%%%%
\section{Position Embeddings: RoPE}
% Source: CS336 Spring 2025, lecture 3 (architecture), slides 30-35
%%%%%%%%%%%%%%%%%%%%%%%%%%%%%%%%%%%%%%%%%%%%%%%%%%%%%%%%%%%%%%%%%%

% FRAMES-BEGIN l3 30 35
% CS336 slide 30
\begin{frame}{Many variations in position embeddings}
% full-body crop (figures would be too small next to the bullets)
\begin{center}
  \includegraphics[page=30,trim=20 3 8 64,clip,width=330pt]{figure/cs336-l3.pdf}
\end{center}
\end{frame}

% CS336 slide 31
\begin{frame}{RoPE: rotary position embeddings}
\begin{center}
  \includegraphics[page=31,trim=62 34 35 67,clip,width=330pt]{figure/cs336-l3.pdf}\par
\end{center}
\end{frame}

% CS336 slide 32
\begin{frame}{RoPE: rotary position embeddings}
% full-body crop (figures would be too small next to the bullets)
\begin{center}
  \includegraphics[page=32,trim=20 3 8 64,clip,width=330pt]{figure/cs336-l3.pdf}
\end{center}
\end{frame}

% CS336 slide 33
\begin{frame}{RoPE: rotary position embeddings}
\begin{itemize}
  \item There are many rotations, which one do you pick?
  \begin{itemize}
    \item Just pair up the coordinates and rotate them in 2d (motivation: complex numbers)
  \end{itemize}
\end{itemize}
\begin{center}
  \includegraphics[page=33,trim=131 53 148 99,clip,width=211pt]{figure/cs336-l3.pdf}\par
\end{center}
{\scriptsize [Su et al 2021]\par}
\end{frame}

% CS336 slide 34
\begin{frame}{The actual RoPE math}
\begin{itemize}
  \item Multiply with sines and cosines
  \item Difference with sine embeddings -- not additive, no cross terms
\end{itemize}
\begin{center}
  \includegraphics[page=34,trim=80 111 72 122,clip,width=330pt]{figure/cs336-l3.pdf}\par
\end{center}
\end{frame}

% CS336 slide 35
\begin{frame}{Implementation and code for RoPE}
\begin{itemize}
  \item \ldots{} Same stuff as the usual multi-head self attention below
  \item Note: embedding at each attention operation to enforce position invariance
\end{itemize}
\begin{center}
  \includegraphics[page=35,trim=46 102 82 91,clip,width=330pt]{figure/cs336-l3.pdf}\par
\end{center}
\end{frame}

% FRAMES-END

%%%%%%%%%%%%%%%%%%%%%%%%%%%%%%%%%%%%%%%%%%%%%%%%%%%%%%%%%%%%%%%%%%
\section{Hyperparameters}
% Source: CS336 Spring 2025, lecture 3 (architecture), slides 36-52
%%%%%%%%%%%%%%%%%%%%%%%%%%%%%%%%%%%%%%%%%%%%%%%%%%%%%%%%%%%%%%%%%%

% FRAMES-BEGIN l3 36 52
% CS336 slide 36
\begin{frame}{Hyperparameters}
\begin{itemize}
  \item Transformer hyperparameter questions you might have had
  \item How much bigger should the feedforward size be compared to hidden size?
  \item How many heads, and should num\_heads always divide hidden size?
  \item What should my vocab size be?
  \item And other model setting questions
  \begin{itemize}
    \item Do people even regularize these huge LMs?
    \item How do people scale these models - very deep or very wide?
  \end{itemize}
\end{itemize}
\end{frame}

% CS336 slide 37
\begin{frame}{Surprising (?) consensus hyperparameter 1}
\begin{center}
  \includegraphics[page=37,trim=62 17 34 67,clip,width=321pt]{figure/cs336-l3.pdf}\par
\end{center}
\end{frame}

% CS336 slide 38
\begin{frame}{Exception \#1 -- GLU variants}
\begin{itemize}
  \item Models are roughly in this range, though PaLM, LLaMA2 and Mistral are slightly larger
\end{itemize}
\begin{center}
  \includegraphics[page=38,trim=62 33 35 67,clip,width=280pt]{figure/cs336-l3.pdf}\par
\end{center}
\end{frame}

% CS336 slide 39
\begin{frame}{Exception \#2 -- T5}
\begin{itemize}
  \item Other, recent exceptions -- Gemma 2 (8x), SmolLM/Gemma 3 (4x, GLU)
\end{itemize}
\begin{center}
  \includegraphics[page=39,trim=62 33 35 67,clip,width=280pt]{figure/cs336-l3.pdf}\par
\end{center}
\end{frame}

% CS336 slide 40
\begin{frame}{Why this range of multipliers?}
\begin{itemize}
  \item Empirically, there's a basin between 1-10 where this hyperparameter is near-optimal
\end{itemize}
\begin{center}
  \includegraphics[page=40,trim=184 62 212 104,clip,width=185pt]{figure/cs336-l3.pdf}\par
\end{center}
{\scriptsize [Kaplan+ 2020]\par}
\end{frame}

% CS336 slide 41
\begin{frame}{What can we learn from the model-dim hyperparam?}
\begin{center}
  \includegraphics[page=41,trim=13 26 39 98,clip,width=330pt]{figure/cs336-l3.pdf}\par
\end{center}
\end{frame}

% CS336 slide 42
\begin{frame}{Surprising (?) consensus hyperparameter 2}
\begin{itemize}
  \item Head-dim*num-heads to model-dim ratio.
  \item This doesn't have to be true: we can have head-dimensions > model-dim / num-heads.
  \begin{itemize}
    \item But most models do follow this guideline
  \end{itemize}
\end{itemize}
\begin{center}
  \includegraphics[page=42,trim=78 142 78 116,clip,width=330pt]{figure/cs336-l3.pdf}\par
\end{center}
\end{frame}

% CS336 slide 43
\begin{frame}{How many heads, whats the model dim?}
\begin{itemize}
  \item Some examples of this hyperparameter
  \item Most models have ratios around 1 -- notable exceptions by some google models.
\end{itemize}
\begin{center}
  \includegraphics[page=43,trim=116 80 62 116,clip,width=313pt]{figure/cs336-l3.pdf}\par
\end{center}
\end{frame}

% CS336 slide 44
\begin{frame}{Evidence for 1-1 ratio?}
\begin{itemize}
  \item There have been papers written against the 1-1 ratio [Bhojanapalli et al 2020]
  \item But we don't seem to be seeing significant `low rank bottlenecks' in practice..
\end{itemize}
\begin{center}
  \includegraphics[page=44,trim=210 95 210 115,clip,width=180pt]{figure/cs336-l3.pdf}\par
\end{center}
\end{frame}

% CS336 slide 45
\begin{frame}{Aspect ratios}
\begin{itemize}
  \item Should my model be deep or wide? How deep and how wide?
  \item Most models are surprisingly consistent on this one too!
\end{itemize}
\begin{center}
  \includegraphics[page=45,trim=16 17 213 146,clip,width=246pt]{figure/cs336-l3.pdf}\par
\end{center}
\end{frame}

% CS336 slide 46
\begin{frame}{Considerations about aspect ratio}
\begin{itemize}
  \item Extremely deep models are harder to parallelize and have higher latency
\end{itemize}
\begin{center}
  \includegraphics[page=46,trim=75 10 26 116,clip,width=304pt]{figure/cs336-l3.pdf}\par
\end{center}
\end{frame}

% CS336 slide 47
\begin{frame}{Evidence on aspect ratio scaling}
% full-body crop (figures would be too small next to the bullets)
\begin{center}
  \includegraphics[page=47,trim=20 3 8 64,clip,width=330pt]{figure/cs336-l3.pdf}
\end{center}
\end{frame}

% CS336 slide 48
\begin{frame}{What are typical vocabulary sizes?}
% full-body crop (figures would be too small next to the bullets)
\begin{center}
  \includegraphics[page=48,trim=20 3 8 64,clip,width=330pt]{figure/cs336-l3.pdf}
\end{center}
\end{frame}

% CS336 slide 49
\begin{frame}{Dropout and other regularization}
\begin{itemize}
  \item Do we need regularization during pretraining?
  \item Arguments against:
  \begin{itemize}
    \item There is a lot of data (trillions of tokens), more than parameters.
    \item SGD only does a single pass on a corpus (hard to memorize)
  \end{itemize}
  \item This is all quite reasonable.. but what do people do in practice?
\end{itemize}
\end{frame}

% CS336 slide 50
\begin{frame}{Dropout and weight decay in practice}
\begin{itemize}
  \item Many older models used dropout during pretraining
  \begin{itemize}
    \item Newer models (except Qwen) rely only on weight decay
  \end{itemize}
\end{itemize}
\begin{center}
  \includegraphics[page=50,trim=181 50 257 80,clip,width=124pt]{figure/cs336-l3.pdf}\par
\end{center}
{\scriptsize * Most of the times papers just don't discuss dropout. On open models, this closely matches not doing dropout. This may not be true of closed models.\par}
\end{frame}

% CS336 slide 51
\begin{frame}{Why weight decay LLMs?}
% full-body crop (figures would be too small next to the bullets)
\begin{center}
  \includegraphics[page=51,trim=20 3 8 64,clip,width=330pt]{figure/cs336-l3.pdf}
\end{center}
\end{frame}

% CS336 slide 52
\begin{frame}{Summary: hyperparameters}
% full-body crop (figures would be too small next to the bullets)
\begin{center}
  \includegraphics[page=52,trim=20 3 8 64,clip,width=330pt]{figure/cs336-l3.pdf}
\end{center}
\end{frame}

% FRAMES-END

%%%%%%%%%%%%%%%%%%%%%%%%%%%%%%%%%%%%%%%%%%%%%%%%%%%%%%%%%%%%%%%%%%
\section{Stability Tricks}
% Source: CS336 Spring 2025, lecture 3 (architecture), slides 53-57
%%%%%%%%%%%%%%%%%%%%%%%%%%%%%%%%%%%%%%%%%%%%%%%%%%%%%%%%%%%%%%%%%%

% FRAMES-BEGIN l3 53 57
% CS336 slide 53
\begin{frame}{Stability tricks}
\begin{itemize}
  \item Recently, lots of attention on stable training
  \begin{itemize}
    \item Don't train models that look like the blue curve!
  \end{itemize}
\end{itemize}
\begin{center}
  \includegraphics[page=53,trim=163 70 163 102,clip,width=205pt]{figure/cs336-l3.pdf}\par
\end{center}
\end{frame}

% CS336 slide 54
\begin{frame}{Where do the issues arise? Beware of softmaxes!}
% full-body crop (figures would be too small next to the bullets)
\begin{center}
  \includegraphics[page=54,trim=20 3 8 64,clip,width=330pt]{figure/cs336-l3.pdf}
\end{center}
\end{frame}

% CS336 slide 55
\begin{frame}{Output softmax stability -- the `z-loss'}
\begin{itemize}
  \item This is useful for stability! PaLM pioneered this `z loss' trick.
  \item Other examples: Baichuan 2 (2023), DCLM (2024), OLMo 2 (2025)
\end{itemize}
\begin{center}
  \includegraphics[page=55,trim=72 192 378 73,clip,width=114pt]{figure/cs336-l3.pdf}\par
  \includegraphics[page=55,trim=366 195 50 99,clip,width=129pt]{figure/cs336-l3.pdf}\par
  \includegraphics[page=55,trim=168 89 132 281,clip,width=178pt]{figure/cs336-l3.pdf}\par
\end{center}
\end{frame}

% CS336 slide 56
\begin{frame}{Attention softmax stability -- the `QK norm'}
\begin{itemize}
  \item Other examples: DCLM, OLMo2, Gemma 2 Originally from vision and multimodal models [Dehgani 2023, Idefcs, Chameleon]
\end{itemize}
\begin{center}
  \includegraphics[page=56,trim=72 85 60 81,clip,width=330pt]{figure/cs336-l3.pdf}\par
\end{center}
\end{frame}

% CS336 slide 57
\begin{frame}{Logit soft-capping.}
\begin{itemize}
  \item Soft-capping the logits to some maximum value via Tanh
  \item Prevents logits from blowing up, but also might have perf issues?
\end{itemize}
\begin{center}
  \includegraphics[page=57,trim=249 141 212 106,clip,width=151pt]{figure/cs336-l3.pdf}\par
  \includegraphics[page=57,trim=202 38 166 319,clip,width=206pt]{figure/cs336-l3.pdf}\par
\end{center}
\end{frame}

% FRAMES-END

%%%%%%%%%%%%%%%%%%%%%%%%%%%%%%%%%%%%%%%%%%%%%%%%%%%%%%%%%%%%%%%%%%
\section{Attention Heads: MQA, GQA, Sliding Windows}
% Source: CS336 Spring 2025, lecture 3 (architecture), slides 58-58
%%%%%%%%%%%%%%%%%%%%%%%%%%%%%%%%%%%%%%%%%%%%%%%%%%%%%%%%%%%%%%%%%%

% FRAMES-BEGIN l3 58 58
% CS336 slide 58
\begin{frame}{Attention heads}
\begin{itemize}
  \item Most models don't touch the attention heads much at all with a few minor exceptions..
  \item GQA / MQA : Saving inference costs by reducing the number of heads
  \item Sparse or sliding window attention (GPT4/Mistral): restricting the attention pattern
  \begin{itemize}
    \item to reduce compute cost
  \end{itemize}
  \item Exotic SSM stuff (Jamba, Falcon 3, etc): not covered (sorry!)
\end{itemize}
\end{frame}

% FRAMES-END

% [O] original bridge frame (draft for the user to rewrite); listed in order.txt
\begin{frame}{The KV Cache}
\begin{itemize}
  \item Autoregressive generation (session 8, decoding) produces one token at a time; each new token attends to all previous tokens.
  \item The keys and values of the previous tokens do not change, so they are stored instead of recomputed: the \textbf{KV cache}.
  \item Its size grows with sequence length $\times$ number of layers $\times$ number of heads $\times$ head dimension (per sequence in the batch).
  \item At inference time, reading this cache from memory --- not the arithmetic --- is often the bottleneck. The following attention variants reduce it.
\end{itemize}
\end{frame}

% FRAMES-BEGIN l3 59 68
% CS336 slide 59
\begin{frame}{GQA/MQA -- Reducing attention head cost}
\begin{itemize}
  \item Let's think about the compute involved for attention
\end{itemize}
\begin{center}
  \includegraphics[page=59,trim=130 127 201 108,clip,width=192pt]{figure/cs336-l3.pdf}\par
  \includegraphics[page=59,trim=63 0 35 298,clip,width=307pt]{figure/cs336-l3.pdf}\par
\end{center}
\end{frame}

% CS336 slide 60
\begin{frame}[shrink=8]{GQA/MQA -- Reducing attention head cost}
\begin{itemize}
  \item What about the incremental case when we generate text?
  \item Key difference: can't parallelize the generation process -- needs to be step by step
  \item In this case -- we need to incrementaly re-compute/update attention via the `KV cache'
\end{itemize}
\begin{center}
  \includegraphics[page=60,trim=191 22 195 194,clip,width=185pt]{figure/cs336-l3.pdf}\par
\end{center}
{\scriptsize [Animation from https://medium.com/@joaolages/kv-caching-explained-276520203249]\par}
\end{frame}

% CS336 slide 61
\begin{frame}{GQA/MQA -- Reducing attention head cost}
\begin{itemize}
  \item What's the incremental arithmetic intensity?
  \item Is there some way around this? The n/d term is difficult to reduce.
\end{itemize}
\begin{center}
  \includegraphics[page=61,trim=62 82 36 140,clip,width=330pt]{figure/cs336-l3.pdf}\par
\end{center}
\end{frame}

% CS336 slide 62
\begin{frame}{MQA -- just have fewer key dimensions.}
\begin{itemize}
  \item Key idea -- have multiple queries, but just one dimension for keys and values
\end{itemize}
\begin{center}
  \includegraphics[page=62,trim=62 1 12 112,clip,width=303pt]{figure/cs336-l3.pdf}\par
\end{center}
\end{frame}

% CS336 slide 63
\begin{frame}{Recent extension -- GQA}
\begin{itemize}
  \item Don't go all the way to one dimension of KV -- have fewer dims
  \item Simple knob to control expressiveness (key-query ratio) and inference efficiency
\end{itemize}
\begin{center}
  \includegraphics[page=63,trim=117 119 130 125,clip,width=284pt]{figure/cs336-l3.pdf}\par
\end{center}
\end{frame}

% CS336 slide 64
\begin{frame}{Does MQA hurt? Sometimes..}
% full-body crop (figures would be too small next to the bullets)
\begin{center}
  \includegraphics[page=64,trim=20 3 8 64,clip,width=330pt]{figure/cs336-l3.pdf}
\end{center}
\end{frame}

% CS336 slide 65
\begin{frame}{Sparse / sliding window attention}
\begin{itemize}
  \item Attending to the entire context can be expensive (quadratic).
  \item Build sparse / structured attention that trades off expressiveness vs runtime (GPT3)
\end{itemize}
\begin{center}
  \includegraphics[page=65,trim=119 39 162 146,clip,width=241pt]{figure/cs336-l3.pdf}\par
\end{center}
{\scriptsize [Child et al 2019]\par}
\end{frame}

% CS336 slide 66
\begin{frame}{Sliding window attention}
\begin{itemize}
  \item Another variation on this idea -- sliding window attention
  \item Just use the main part of the strided pattern -- let depth extend effective context (Mistral)
\end{itemize}
\begin{center}
  \includegraphics[page=66,trim=90 93 118 106,clip,width=301pt]{figure/cs336-l3.pdf}\par
\end{center}
\end{frame}

% CS336 slide 67
\begin{frame}{Current standard trick -- interleave `full' and `LR' attention}
\begin{itemize}
  \item From Cohere Command A -- Every 4th layer is a full attention
  \item Long-range info via NoPE, short-range info via RoPE + SWA.
  \item Other models -- LLaMA 4, Gemma does SWA+Full RoPE.
\end{itemize}
\begin{center}
  \includegraphics[page=67,trim=178 104 178 104,clip,width=193pt]{figure/cs336-l3.pdf}\par
\end{center}
\end{frame}

% CS336 slide 68
\begin{frame}{Recap, conclusion, etc.}
\begin{itemize}
  \item Many aspects (arch, hparams) of transformers are in common across the big LMs
  \begin{itemize}
    \item Major differences? Position embeddings, activations, tokenization
  \end{itemize}
\end{itemize}
\begin{center}
  \includegraphics[page=68,trim=35 92 21 106,clip,width=330pt]{figure/cs336-l3.pdf}\par
\end{center}
\end{frame}

% FRAMES-END

%%%%%%%%%%%%%%%%%%%%%%%%%%%%%%%%%%%%%%%%%%%%%%%%%%%%%%%%%%%%%%%%%%
\section{Mixture of Experts}
% Source: CS336 Spring 2025, lecture 4 (MoEs), slides 2-12
%%%%%%%%%%%%%%%%%%%%%%%%%%%%%%%%%%%%%%%%%%%%%%%%%%%%%%%%%%%%%%%%%%

% FRAMES-BEGIN l4 2 12
% CS336 slide 2
\begin{frame}{Mixture of experts}
\begin{center}
  \includegraphics[page=2,trim=30 24 12 68,clip,width=330pt]{figure/cs336-l4.pdf}\par
\end{center}
\end{frame}

% CS336 slide 3
\begin{frame}{What's a MoE?}
\begin{itemize}
  \item Replace big feedforward with (many) big feedforward networks and a selector layer
  \begin{itemize}
    \item You can increase the \# experts without affecting FLOPs
  \end{itemize}
\end{itemize}
\begin{center}
  \includegraphics[page=3,trim=50 108 34 100,clip,width=330pt]{figure/cs336-l4.pdf}\par
\end{center}
\end{frame}

% CS336 slide 4
\begin{frame}{Why are MoEs getting popular?}
\begin{itemize}
  \item Same FLOP, more param does better
\end{itemize}
\begin{center}
  \includegraphics[page=4,trim=79 29 42 112,clip,width=310pt]{figure/cs336-l4.pdf}\par
\end{center}
\end{frame}

% CS336 slide 5
\begin{frame}{Why are MoEs getting popular?}
\begin{center}
  \includegraphics[page=5,trim=62 15 6 84,clip,width=330pt]{figure/cs336-l4.pdf}\par
\end{center}
\end{frame}

% CS336 slide 6
\begin{frame}{Why are MoEs getting popular?}
\begin{itemize}
  \item Highly competitive vs dense equivalents
\end{itemize}
\begin{center}
  \includegraphics[page=6,trim=183 55 183 88,clip,width=185pt]{figure/cs336-l4.pdf}\par
\end{center}
\end{frame}

% CS336 slide 7
\begin{frame}{Why are MoEs getting popular?}
\begin{center}
  \includegraphics[page=7,trim=72 14 200 72,clip,width=231pt]{figure/cs336-l4.pdf}\par
\end{center}
\end{frame}

% CS336 slide 8
\begin{frame}{Some MoE results -- from the west}
\begin{itemize}
  \item MoEs are most of the highest-performance open models, and are quite quick
\end{itemize}
\begin{center}
  \includegraphics[page=8,trim=116 82 116 82,clip,width=277pt]{figure/cs336-l4.pdf}\par
\end{center}
\end{frame}

% CS336 slide 9
\begin{frame}{Earlier MoE results from Chinese groups -- Qwen}
\begin{itemize}
  \item Chinese LLM companies are also doing quite a bit of MoE work on the smaller end
\end{itemize}
\begin{center}
  \includegraphics[page=9,trim=71 106 315 136,clip,width=140pt]{figure/cs336-l4.pdf}\par
  \includegraphics[page=9,trim=408 106 54 136,clip,width=109pt]{figure/cs336-l4.pdf}\par
\end{center}
\end{frame}

% CS336 slide 10
\begin{frame}{Earlier MoE results from Chinese groups - DeepSeek}
\begin{itemize}
  \item There's also some good recent ablation work on MoEs showing they're generally good
\end{itemize}
\begin{center}
  \includegraphics[page=10,trim=193 20 193 107,clip,width=165pt]{figure/cs336-l4.pdf}\par
\end{center}
\end{frame}

% CS336 slide 11
\begin{frame}{Recent MoE results -- DeepSeek v3}
\begin{center}
  \includegraphics[page=11,trim=119 30 82 67,clip,width=277pt]{figure/cs336-l4.pdf}\par
\end{center}
\end{frame}

% CS336 slide 12
\begin{frame}{Why haven't MoEs been more popular?}
\begin{itemize}
  \item Infrastructure is complex / advantages on multi node
  \item Training objectives are somewhat heuristic (and sometimes unstable)
\end{itemize}
\begin{center}
  \includegraphics[page=12,trim=112 166 34 103,clip,width=249pt]{figure/cs336-l4.pdf}\par
  \includegraphics[page=12,trim=193 -0 220 262,clip,width=133pt]{figure/cs336-l4.pdf}\par
\end{center}
{\scriptsize [Zoph et al 2022]\par}
\end{frame}

% FRAMES-END

%%%%%%%%%%%%%%%%%%%%%%%%%%%%%%%%%%%%%%%%%%%%%%%%%%%%%%%%%%%%%%%%%%
\section{MoE Routing}
% Source: CS336 Spring 2025, lecture 4 (MoEs), slides 13-23
%%%%%%%%%%%%%%%%%%%%%%%%%%%%%%%%%%%%%%%%%%%%%%%%%%%%%%%%%%%%%%%%%%

% FRAMES-BEGIN l4 13 23
% CS336 slide 13
\begin{frame}{What MoEs generally look like}
% full-body crop (figures would be too small next to the bullets)
\begin{center}
  \includegraphics[page=13,trim=20 3 8 64,clip,width=330pt]{figure/cs336-l4.pdf}
\end{center}
\end{frame}

% CS336 slide 14
\begin{frame}{MoE -- what varies?}
\begin{itemize}
  \item Routing function
  \item Expert sizes
  \item Training objectives
\end{itemize}
\end{frame}

% CS336 slide 15
\begin{frame}{Routing function - overview}
\begin{itemize}
  \item Many of the routing algorithms boil down to `choose top k'
\end{itemize}
\begin{center}
  \includegraphics[page=15,trim=122 19 28 118,clip,width=291pt]{figure/cs336-l4.pdf}\par
\end{center}
\end{frame}

% CS336 slide 16
\begin{frame}{Routing type}
\begin{itemize}
  \item Almost all the MoEs do a standard `token choice topk' routing. Some recent ablations
\end{itemize}
\begin{center}
  \includegraphics[page=16,trim=51 119 51 119,clip,width=330pt]{figure/cs336-l4.pdf}\par
\end{center}
\end{frame}

% CS336 slide 17
\begin{frame}{Common routing variants in detail}
% full-body crop (figures would be too small next to the bullets)
\begin{center}
  \includegraphics[page=17,trim=20 3 8 64,clip,width=330pt]{figure/cs336-l4.pdf}
\end{center}
\end{frame}

% CS336 slide 18
\begin{frame}{Other routing methods}
\begin{center}
  \includegraphics[page=18,trim=12 176 26 71,clip,width=330pt]{figure/cs336-l4.pdf}\par
  \includegraphics[page=18,trim=18 8 6 240,clip,width=330pt]{figure/cs336-l4.pdf}\par
\end{center}
\end{frame}

% CS336 slide 19
\begin{frame}{Top-K routing in detail.}
\begin{itemize}
  \item Most papers do the old and classic top-k routing. How does this work?
  \begin{itemize}
    \item This is the DeepSeek (V1-2) router (Grok, Qwen do this too)
    \item Mixtral, DBRX, DeepSeek v3 softmaxes after the TopK
  \end{itemize}
  \item Gates selected by a logistic regressor
\end{itemize}
\begin{center}
  \includegraphics[page=19,trim=21 46 275 112,clip,width=153pt]{figure/cs336-l4.pdf}\par
\end{center}
{\scriptsize [Dai et al 2024]\par}
\end{frame}

% CS336 slide 20
\begin{frame}{Recent variations from DeepSeek and other Chinese LMs}
\begin{itemize}
  \item Smaller, larger number of experts + a few shared experts that are always on.
  \begin{itemize}
    \item (Used in DeepSeek / Qwen, originally from DeepSpeed MoE)
  \end{itemize}
\end{itemize}
\begin{center}
  \includegraphics[page=20,trim=117 86 81 64,clip,width=247pt]{figure/cs336-l4.pdf}\par
\end{center}
\end{frame}

% CS336 slide 21
\begin{frame}{Various ablations from the DeepSeek paper}
\begin{itemize}
  \item More experts, shared experts all seem to generally help
\end{itemize}
\begin{center}
  \includegraphics[page=21,trim=128 66 128 69,clip,width=236pt]{figure/cs336-l4.pdf}\par
\end{center}
\end{frame}

% CS336 slide 22
\begin{frame}{Ablations from OlMoE}
\begin{itemize}
  \item Gains from fine-grained experts, none from shared experts.
\end{itemize}
\begin{center}
  \includegraphics[page=22,trim=70 4 80 100,clip,width=260pt]{figure/cs336-l4.pdf}\par
\end{center}
\end{frame}

% CS336 slide 23
\begin{frame}{Expert routing setups for recent MoEs}
\begin{center}
  \includegraphics[page=23,trim=63 20 71 63,clip,width=300pt]{figure/cs336-l4.pdf}\par
\end{center}
\end{frame}

% FRAMES-END

%%%%%%%%%%%%%%%%%%%%%%%%%%%%%%%%%%%%%%%%%%%%%%%%%%%%%%%%%%%%%%%%%%
\section{Training MoEs}
% Source: CS336 Spring 2025, lecture 4 (MoEs), slides 24-37
%%%%%%%%%%%%%%%%%%%%%%%%%%%%%%%%%%%%%%%%%%%%%%%%%%%%%%%%%%%%%%%%%%

% FRAMES-BEGIN l4 24 37
% CS336 slide 24
\begin{frame}{How do we train MoEs?}
\begin{itemize}
  \item Major challenge: we need sparsity for training-time efficiency\ldots{}
  \begin{itemize}
    \item But sparse gating decisions are not differentiable!
  \end{itemize}
  \item Solutions?
  \item 1. Reinforcment learning to optimize gating policies 2. Stochastic perturbations 3. Heuristic `balancing' losses.
  \begin{itemize}
    \item Guess which one people use in practice?
  \end{itemize}
\end{itemize}
\end{frame}

% CS336 slide 25
\begin{frame}{RL for MoEs}
\begin{itemize}
  \item RL via REINFORCE does work, but not so much better that it's a clear win
  \begin{itemize}
    \item RL is the `right solution' but gradient variances and complexity
    \item means it's not widely used
  \end{itemize}
\end{itemize}
\begin{center}
  \includegraphics[page=25,trim=55 122 35 137,clip,width=330pt]{figure/cs336-l4.pdf}\par
\end{center}
{\scriptsize (REINFORCE baseline approach, Clark et al 2020)\par}
\end{frame}

% CS336 slide 26
\begin{frame}{Stochastic approximations}
\begin{itemize}
  \item From Shazeer et al 2017 -- routing decisions are stochastic with gaussian perturbations.
  \item 1. This naturally leads to experts that are a bit more robust. 2. The softmax means that the model learns how to rank K experts
\end{itemize}
\begin{center}
  \includegraphics[page=26,trim=106 137 106 90,clip,width=305pt]{figure/cs336-l4.pdf}\par
\end{center}
\end{frame}

% CS336 slide 27
\begin{frame}{Stochastic approximations}
\begin{itemize}
  \item Stochastic jitter in Fedus et al 2022. This does a uniform multiplicative perturbation for the
  \item same goal of getting less brittle experts. This was later removed in Zoph et al 2022
\end{itemize}
\begin{center}
  \includegraphics[page=27,trim=98 187 61 86,clip,width=315pt]{figure/cs336-l4.pdf}\par
  \includegraphics[page=27,trim=212 21 212 301,clip,width=167pt]{figure/cs336-l4.pdf}\par
\end{center}
\end{frame}

% CS336 slide 28
\begin{frame}{Heuristic balancing losses}
\begin{itemize}
  \item Another key issue -- systems efficiency requires that we use experts evenly..
\end{itemize}
\begin{center}
  \includegraphics[page=28,trim=21 8 13 102,clip,width=319pt]{figure/cs336-l4.pdf}\par
\end{center}
\end{frame}

% CS336 slide 29
\begin{frame}{Example from deepseek (v1-2)}
\begin{itemize}
  \item Per-expert balancing -- same as the switch transformer
  \item Per-device balancing -- the objective above, but aggregated by device.
\end{itemize}
\begin{center}
  \includegraphics[page=29,trim=114 182 186 109,clip,width=230pt]{figure/cs336-l4.pdf}\par
  \includegraphics[page=29,trim=111 17 183 281,clip,width=233pt]{figure/cs336-l4.pdf}\par
\end{center}
\end{frame}

% CS336 slide 30
\begin{frame}{DeepSeek v3 variation -- per-expert biases}
\begin{itemize}
  \item Set up a per-expert bias (making it more likely to get tokens) and use online learning
  \item They call this `auxiliary loss free balancing'
\end{itemize}
\begin{center}
  \includegraphics[page=30,trim=172 237 172 104,clip,width=171pt]{figure/cs336-l4.pdf}\par
  \includegraphics[page=30,trim=72 6 142 198,clip,width=231pt]{figure/cs336-l4.pdf}\par
\end{center}
\end{frame}

% CS336 slide 31
\begin{frame}{What happens when removing load balancing losses?}
\begin{center}
  \includegraphics[page=31,trim=158 14 158 74,clip,width=210pt]{figure/cs336-l4.pdf}\par
\end{center}
\end{frame}

% CS336 slide 32
\begin{frame}{Training MoEs -- the systems side}
% full-body crop (figures would be too small next to the bullets)
\begin{center}
  \includegraphics[page=32,trim=20 3 8 64,clip,width=330pt]{figure/cs336-l4.pdf}
\end{center}
\end{frame}

% CS336 slide 33
\begin{frame}{Training MoEs -- the systems side}
\begin{itemize}
  \item MoE routing allows for parallelism, but also some complexities
  \item Modern libraries like MegaBlocks (used in many open MoEs) use smarter sparse MMs
\end{itemize}
\begin{center}
  \includegraphics[page=33,trim=56 112 51 125,clip,width=330pt]{figure/cs336-l4.pdf}\par
\end{center}
\end{frame}

% CS336 slide 34
\begin{frame}{Fun side issue -- stochasticity of MoE models}
\begin{itemize}
  \item There was speculation that GPT-4's stochasticity was due to MoE..
  \item Why would a MoE have additional randomness?
  \begin{itemize}
    \item Token dropping from routing happens at a batch level -- this means that
    \item other people's queries can drop your token!
  \end{itemize}
\end{itemize}
\begin{center}
  \includegraphics[page=34,trim=101 79 101 155,clip,width=270pt]{figure/cs336-l4.pdf}\par
\end{center}
\end{frame}

% CS336 slide 35
\begin{frame}{Issues with MoEs - stability}
\begin{itemize}
  \item Solution: Use Float 32 just for the expert router (sometimes with an aux z-loss)
\end{itemize}
\begin{center}
  \includegraphics[page=35,trim=54 213 507 76,clip,width=70pt]{figure/cs336-l4.pdf}\par
  \includegraphics[page=35,trim=220 194 20 84,clip,width=210pt]{figure/cs336-l4.pdf}\par
  \includegraphics[page=35,trim=72 64 143 271,clip,width=221pt]{figure/cs336-l4.pdf}\par
\end{center}
\end{frame}

% CS336 slide 36
\begin{frame}{Z-loss stability for the router}
\begin{itemize}
  \item What happens when we remove the z-loss?
\end{itemize}
\begin{center}
  \includegraphics[page=36,trim=88 100 88 100,clip,width=327pt]{figure/cs336-l4.pdf}\par
\end{center}
\end{frame}

% CS336 slide 37
\begin{frame}{Issues with MoEs -- fine-tuning}
% full-body crop (figures would be too small next to the bullets)
\begin{center}
  \includegraphics[page=37,trim=20 3 8 64,clip,width=330pt]{figure/cs336-l4.pdf}
\end{center}
\end{frame}

% FRAMES-END

%%%%%%%%%%%%%%%%%%%%%%%%%%%%%%%%%%%%%%%%%%%%%%%%%%%%%%%%%%%%%%%%%%
\section{Upcycling and DeepSeek MoE}
% Source: CS336 Spring 2025, lecture 4 (MoEs), slides 38-47
%%%%%%%%%%%%%%%%%%%%%%%%%%%%%%%%%%%%%%%%%%%%%%%%%%%%%%%%%%%%%%%%%%

% FRAMES-BEGIN l4 38 47
% CS336 slide 38
\begin{frame}{Other training methods - upcycling}
\begin{itemize}
  \item Can we use a pre-trained LM to initialize a MoE?
\end{itemize}
\begin{center}
  \includegraphics[page=38,trim=459 110 30 101,clip,width=84pt]{figure/cs336-l4.pdf}\par
  \includegraphics[page=38,trim=33 117 271 104,clip,width=150pt]{figure/cs336-l4.pdf}\par
\end{center}
\end{frame}

% CS336 slide 39
\begin{frame}{Upcycling example - MiniCPM}
\begin{itemize}
  \item Uses the MiniCPM model (topk=2, 8 experts, \textasciitilde{} 4B active params).
  \item Simple MoE, shows gains from the base model with \textasciitilde{} 520B tokens for training
\end{itemize}
\begin{center}
  \includegraphics[page=39,trim=148 132 115 126,clip,width=274pt]{figure/cs336-l4.pdf}\par
\end{center}
\end{frame}

% CS336 slide 40
\begin{frame}{Upcycling example -- Qwen MoE}
\begin{itemize}
  \item Qwen MoE -- Initialized from the Qwen 1.8B model top-k=4, 60 experts w/ 4 shared.
  \item Similar architecture / setup to DeepSeekMoE, but one of the first (confirmed) upcycling successes
\end{itemize}
\begin{center}
  \includegraphics[page=40,trim=117 114 109 135,clip,width=297pt]{figure/cs336-l4.pdf}\par
\end{center}
\end{frame}

% CS336 slide 41
\begin{frame}{DeepSeek MoE v1-v2-v3}
\begin{itemize}
  \item To wrap up, we'll walk through the DeepSeek MoE architecture.
\end{itemize}
\begin{center}
  \includegraphics[page=41,trim=78 149 217 90,clip,width=175pt]{figure/cs336-l4.pdf}\par
  \includegraphics[page=41,trim=90 40 466 272,clip,width=67pt]{figure/cs336-l4.pdf}\par
  \includegraphics[page=41,trim=386 34 194 296,clip,width=58pt]{figure/cs336-l4.pdf}\par
\end{center}
{\scriptsize Standard Aux-loss balancing (Expert + Device)\par}
\end{frame}

% CS336 slide 42
\begin{frame}{DeepSeek MoE v2}
\begin{center}
  \includegraphics[page=42,trim=70 174 212 69,clip,width=193pt]{figure/cs336-l4.pdf}\par
  \includegraphics[page=42,trim=58 54 414 246,clip,width=109pt]{figure/cs336-l4.pdf}\par
  \includegraphics[page=42,trim=374 4 108 293,clip,width=105pt]{figure/cs336-l4.pdf}\par
\end{center}
{\scriptsize Communication balancing loss -- balancing both communication in and out\par}
\end{frame}

% CS336 slide 43
\begin{frame}{DeepSeek MoE v3}
\begin{center}
  \includegraphics[page=43,trim=70 12 80 69,clip,width=291pt]{figure/cs336-l4.pdf}\par
\end{center}
\end{frame}

% CS336 slide 44
\begin{frame}{Bonus: What else do you need to make DeepSeek MoE v3?}
\begin{itemize}
  \item MLA : Multihead, latent attention
  \item Basic idea: express the Q, K, V as functions of a lower-dim, `latent' activation
\end{itemize}
\begin{center}
  \includegraphics[page=44,trim=134 86 134 103,clip,width=254pt]{figure/cs336-l4.pdf}\par
\end{center}
\end{frame}

% CS336 slide 45
\begin{frame}{What else do you need to make DeepSeek MoE v3?}
\begin{itemize}
  \item Basic idea: express the Q, K, V as functions of a lower-dim, `latent' activation
\end{itemize}
\begin{center}
  \includegraphics[page=45,trim=62 2 35 95,clip,width=277pt]{figure/cs336-l4.pdf}\par
\end{center}
\end{frame}

% CS336 slide 46
\begin{frame}{What else do you need to make DeepSeek MoE v3?}
\begin{itemize}
  \item MTP: Have small, lightweight models that predict multiple steps ahead
  \begin{itemize}
    \item (See paper for ablations)
  \end{itemize}
\end{itemize}
\begin{center}
  \includegraphics[page=46,trim=39 90 15 98,clip,width=272pt]{figure/cs336-l4.pdf}\par
  \includegraphics[page=46,trim=224 4 213 322,clip,width=116pt]{figure/cs336-l4.pdf}\par
\end{center}
\end{frame}

% CS336 slide 47
\begin{frame}{MoE summary}
\begin{itemize}
  \item MoEs take advantage of sparsity -- not all inputs need the full model
  \begin{itemize}
    \item Discrete routing is hard, but top-k heuristics seem to work
  \end{itemize}
  \item Lots of empirical evidence now that MoEs work, and are cost-effective
\end{itemize}
\end{frame}

% FRAMES-END

%%%%%%%%%%%%%%%%%%%%%%%%%%%%%%%%%%%%%%%%%%%%%%%%%%%%%%%%%%%%%%%%%%
\section{Where This Is Going}
%%%%%%%%%%%%%%%%%%%%%%%%%%%%%%%%%%%%%%%%%%%%%%%%%%%%%%%%%%%%%%%%%%

% [O] original bridge frame (draft for the user to rewrite); listed in order.txt
\begin{frame}{Where This Is Going}
\begin{itemize}
  \item The changes in this session act on the Transformer \emph{block}, not on the type of data.
  \item Vision Transformers use them too, e.g.:
  \begin{itemize}
    \item ViT-22B: parallel layers, QK normalization, no biases on the QKV projections and LayerNorms \citebutton{Dehghani et al., 2023}{https://arxiv.org/abs/2302.05442};
    \item V-MoE: a sparse mixture-of-experts Vision Transformer \citebutton{Riquelme et al., 2021}{https://arxiv.org/abs/2106.05974}.
  \end{itemize}
  \item Session 10 takes the Transformer to images, multimodal data and video.
\end{itemize}
\end{frame}


%%%%%%%%%%%%%%%%%%%%%%%%%%%%%%%%%%%%%%%%%%%%%%%%%%%%%%%%%%%%%%%%%%
\section{License \& Attribution}
%%%%%%%%%%%%%%%%%%%%%%%%%%%%%%%%%%%%%%%%%%%%%%%%%%%%%%%%%%%%%%%%%%

\begin{frame}{License \& Attribution}
\begin{itemize}
\item \textbf{Source.} All frames except the four listed below are adapted from \textbf{Stanford CS336 \emph{Language Modeling from Scratch}}, Spring 2025, Lectures 3 (\emph{Architectures, hyperparameters}) and 4 (\emph{Mixture of experts}) by Tatsunori Hashimoto, \url{https://github.com/stanford-cs336/spring2025-lectures} --- \textbf{Copyright 2025 Stanford University}, licensed under the \textbf{MIT License} (full notice in \texttt{\_sources/}). One frame per source slide, in the source's order; text in the source's wording, figures and formulas cropped from the source PDF.
\item \textbf{Figures} inside the slides are from the papers credited on each frame, reused with that attribution.
\item \textbf{Original to this course.} \enquote{Since 2017: What Changed Inside the Transformer?}, \enquote{Compute vs.\ Memory}, \enquote{The KV Cache}, \enquote{Where This Is Going}.
\end{itemize}
\end{frame}

\endlecture
\end{document}
